\ELhold

\ELsection{سرفصل‌ها}{سرفصل‌ها}

\begin{ELunit}

\ELfigure{m05-course-outline}{}

\end{ELunit}

\ELhold

\ELsection{میدان‌های مغناطیسی ساکن}{میدان‌های مغناطیسی ساکن}

\begin{ELunit}

\begin{itemize}
\tightlist
\item
  هنگامیکه بار آزمون کوچکی در میدان الکتریکی \ELim{\vect{E}} قرار می‌گیرد، نیروی زیر به آن وارد می‌شود
\end{itemize}

\ELdisp{\vect{F}_e\ELbr =q\vect{E}\qquad\ELbr (\mathrm{N})}

\begin{itemize}
\tightlist
\item
  هنگامیکه این بار آزمون، در حال حرکت درون یک میدان مغناطیسی باشد، نیروی دیگری به آن وارد می‌شود که برابر است با
\end{itemize}

\ELdisp{\vect{F}_m\ELbr =q\vect{u}\times\vect{B}\qquad\ELbr (\mathrm{N})}

\end{ELunit}

\begin{ELjoin}

\begin{itemize}
\tightlist
\item
  \ELim{\vect{B}}: چگالی شار مغناطیسی با واحد وبر بر متر مربع یا تسلا
\item
  پس نیروی الکترومغناطیسی کل وارد بر بار \ELim{q} برابر است با
\end{itemize}

\begin{ELimportant}{معادله نیروی لورنتس}

\ELdisp{\vect{F}\ELbr =q(\vect{E}+\vect{u}\times\vect{B})\qquad\ELbr (\mathrm{N})}

\end{ELimportant}

\end{ELjoin}

\begin{ELunit}

\ELfigure{m05-map}{}

\begin{itemize}
\tightlist
\item
  فرضیات اصلی مغناطیس ساکن که بر اساس آن‌ها می‌توان دیگر روابط و قوانین این حوزه را استخراج کرد کدامند؟
\item
  چگونه بر اساس فرضیات فوق، \textbf{قانون مداری آمپر} استخراج می‌شود؟
\item
  چگونه می‌توان با استفاده از قانون مداری آمپر، چگالی شار مغناطیسی ناشی از یک جریان را بدست آورد؟
\end{itemize}

\end{ELunit}

\ELhold

\ELsection{فرضیات اصلی مغناطیس ساکن در فضای آزاد}{فرضیات اصلی مغناطیس ساکن در فضای آزاد}

\begin{ELjoin}

\begin{itemize}
\tightlist
\item
  دو فرض اساسی مغناطیس ساکن در فضای آزاد
\end{itemize}

\begin{ELimportant}{}

\ELdisp{\nabla\cdot\vect{B}\ELbr =0}
\ELdisp{\nabla\times\vect{B}\ELbr =\mu_0\vect{J}}

\end{ELimportant}

\end{ELjoin}

\begin{ELunit}

\begin{itemize}
\tightlist
\item
  \ELim{\mu_0=4\pi\times10^{-7}\,\mathrm{(H/m)}}: ضریب نفوذپذیری فضای آزاد
\item
  \ELim{\vect{J}}: چگالی حجمی جریان \ELim{\mathrm{(A/m^2)}}
\item
  چون دیورژانس کرل هر میدان برداری صفر است، در نتیجه
\end{itemize}

\ELdisp{\nabla\cdot\vect{J}\ELbr =0}

\end{ELunit}

\begin{ELjoin}

\begin{itemize}
\tightlist
\item
  برای چگالی بار الکتریکی هیچ مشابه مغناطیسی وجود ندارد.
\item
  شکل انتگرالی رابطه اول \ELim{\ELto} انتگرالگیری روی حجم دلخواه \ELim{V} و استفاده از قضیه دیورژانس
\end{itemize}

\begin{ELimportant}{قانون بقای شار مغناطیسی}

\ELdisp{\oint_S\vect{B}\cdot d\vect{s}\ELbr =0}

\end{ELimportant}

\end{ELjoin}

\begin{ELunit}

\begin{itemize}
\tightlist
\item
  این معادله را \textbf{قانون بقای شار مغناطیسی} می‌نامند.

  \begin{itemize}
  \tightlist
  \item
    هیچ منبع شار مغناطیسی وجود ندارد و خطوط شار مغناطیسی همیشه در خود بسته می‌شوند.
  \end{itemize}
\item
  شکل انتگرالی رابطه دوم \ELim{\ELto} انتگرالگیری روی سطح باز دلخواه \ELim{S} و استفاده از قضیه استوکس
\end{itemize}

\ELdisp{\int_S(\nabla\times\vect{B})\cdot d\vect{s}\ELbr =\mu_0\int_S\vect{J}\cdot d\vect{s}}

\end{ELunit}

\begin{ELimportant}{قانون مداری آمپر}

\ELdisp{\oint_C\vect{B}\cdot\dif\vect{\ell}\ELbr =\mu_0I}

\end{ELimportant}

\begin{ELunit}

\begin{itemize}
\tightlist
\item
  \ELim{C}: مسیر محصور کننده سطح \ELim{S} (مسیر \ELim{C} و جریان \ELim{I} از قانون دست راست تبعیت می‌کنند)
\item
  \ELim{I}: جریان گذرنده از سطح \ELim{S}
\item
  این معادله شکلی از \textbf{قانون مداری آمپر} است.

  \begin{itemize}
  \tightlist
  \item
    گردش چگالی شار مغناطیسی در فضای آزاد به دور هر مسیر بسته، برابر با حاصلضرب \ELim{\mu_0} در کل جریان گذرنده از سطح محصور شده توسط این مسیر است.
  \end{itemize}
\item
  قانون مداری آمپر در تعیین چگالی شار مغناطیسی ناشی از جریان \ELim{I} هنگامیکه مسیر بسته \ELim{C} به دور جریان چنان وجود داشته باشد که \ELim{\vect{B}} روی مسیر ثابت باشد، مفید است.
\end{itemize}

\end{ELunit}

\ELnextnum{5-1}

\begin{ELexample}{}

یک هادی بسیار بلند مستقیم با سطح مقطع دایروی به شعاع \ELim{b} حامل جریان \ELim{I} است. چگالی شار مغناطیسی درون و بیرون هادی را بدست آورید.

\ELfigure{m05-ex1}{}

\begin{ELsolution}

\begin{itemize}
\tightlist
\item
  اگر هادی در امتداد محور \ELim{z} باشد، \ELim{\vect{B}} در راستای \ELim{\phi} بوده و مقدار آن در امتداد هر مسیر دایروی حول محور \ELim{z} ثابت است.

  \begin{itemize}
  \tightlist
  \item
    درون هادی
  \end{itemize}
\end{itemize}

\ELdisp{\left.\begin{aligned}&\vect{B}_1=\uvec{\phi}B_{\phi1},\qquad\dif\vect{\ell}=\uvec{\phi}r_1\,d\phi\\&\oint_{C_1}\vect{B}_1\cdot\dif\vect{\ell}=\int_0^{2\pi}B_{\phi1}r_1\,d\phi=2\pi r_1B_{\phi1}\\&I_1=\frac{\pi r_1^2}{\pi b^2}I=\left(\frac{r_1}{b}\right)^2I\end{aligned}\right\}\quad\ELbr \vect{B}_1\ELbr =\uvec{\phi}B_{\phi1}\ELbr =\uvec{\phi}\frac{\mu_0r_1I}{2\pi b^2}}

\begin{itemize}
\tightlist
\item
  بیرون هادی
\end{itemize}

\ELdisp{\vect{B}_2\ELbr =\uvec{\phi}B_{\phi2},\qquad\ELbr \dif\vect{\ell}\ELbr =\uvec{\phi}r_2\,d\phi}
\ELdisp{\oint_{C_2}\vect{B}_2\cdot\dif\vect{\ell}\ELbr =2\pi r_2B_{\phi2}}
\ELdisp{\vect{B}_2\ELbr =\uvec{\phi}B_{\phi2}\ELbr =\uvec{\phi}\frac{\mu_0I}{2\pi r_2}}

\end{ELsolution}

\end{ELexample}

\ELnextnum{5-2}

\begin{ELexample}{}

چگالی شار مغناطیسی درون یک سیم پیچ چنبره‌ای با هسته هوا و \ELim{N} دور سیمِ حامل جریان \ELim{I} را بدست آورید.

\ELfigure{m05-ex2}{}

\begin{ELsolution}

\begin{itemize}
\tightlist
\item
  تقارن استوانه‌ای تضمین می‌کند که \ELim{\vect{B}} فقط مولفه در راستای \ELim{\phi} دارد و مقدار آن در امتداد هر مسیر دایروی حول محور سیم پیچ ثابت است.
\end{itemize}

\ELdisp{\oint\vect{B}\cdot\dif\vect{\ell}\ELbr =2\pi rB_\phi\ELbr =\mu_0NI}
\ELdisp{\vect{B}\ELbr =\uvec{\phi}B_\phi\ELbr =\uvec{\phi}\frac{\mu_0NI}{2\pi r},\qquad\ELbr (b-a)<r<(b+a)}

\end{ELsolution}

\end{ELexample}

\ELnextnum{5-3}

\begin{ELexample}{}

چگالی شار مغناطیسی درون یک سیم پیچ بسیار بلند با هسته هوا و \ELim{n} دور سیم در واحد طول حامل جریان \ELim{I} را بدست آورید.

\ELfigure{m05-ex3}{}

\begin{ELsolution}

\begin{itemize}
\tightlist
\item
  خارج از سیم پیچ میدان مغناطیسی صفر است.
\item
  با توجه به تقارن ساختار، \ELim{\vect{B}} هم‌راستا با محور سیم پیچ بوده و مقدار آن در امتداد سیم پیچ ثابت است.
\end{itemize}

\ELdisp{BL\ELbr =\mu_0nLI}
\ELdisp{B\ELbr =\mu_0nI}

\end{ELsolution}

\end{ELexample}

\ELhold

\ELsection{پتانسیل مغناطیسی برداری}{پتانسیل مغناطیسی برداری}

\begin{ELunit}

\ELfigure{m05-map-vecpot}{}

\begin{itemize}
\tightlist
\item
  مفهوم پتانسیل مغناطیسی برداری چیست، کاربرد آن کجاست و چگونه محاسبه می‌شود؟
\item
  رابطه قانون بیوساوار و نحوه استفاده از آن برای محاسبه چگالی شار مغناطیسی چگونه است؟
\end{itemize}

\ELdisp{\nabla\cdot\vect{B}\ELbr =0\quad\ELbr \ELbr \Longrightarrow\quad\ELbr \boxed{\vect{B}=\nabla\times\vect{A}\qquad(\mathrm{T})}}

\begin{itemize}
\tightlist
\item
  \ELim{\vect{A}} را \textbf{پتانسیل مغناطیسی برداری} می‌نامند. واحد آن وبر بر متر است.

  \begin{itemize}
  \tightlist
  \item
    اگر بتوانیم در مورد یک توزیع جریان، \ELim{\vect{A}} را بیابیم، آنگاه به سادگی \ELim{\vect{B}} قابل محاسبه خواهد بود (دقیقاً مانند \ELim{V} و \ELim{\vect{E}}).
  \end{itemize}
\item
  تعریف یک بردار به مشخص کردن کرل و دیورژانس آن نیاز دارد. برای مشخص کردن دیورژانس \ELim{\vect{A}} به ترتیب زیر عمل می‌کنیم
\end{itemize}

\ELdisp{\nabla\times\vect{B}\ELbr =\mu_0\vect{J}\quad\ELbr \ELbr \Longrightarrow\quad\ELbr \nabla\times\nabla\times\vect{A}\ELbr =\mu_0\vect{J}}

\begin{itemize}
\tightlist
\item
  \textbf{لاپلاسین یک کمیت برداری} را به صورت زیر تعریف می‌کنیم
\end{itemize}

\ELdisp{\nabla^2\vect{A}\ELbr =\nabla(\nabla\cdot\vect{A})-\nabla\times\nabla\times\vect{A}}

\begin{itemize}
\tightlist
\item
  به عنوان مثال لاپلاسین \ELim{\vect{A}} در دستگاه مختصات کارتزین به صورت زیر است
\end{itemize}

\ELdisp{\nabla^2\vect{A}\ELbr =\uvec{x}\nabla^2A_x+\uvec{y}\nabla^2A_y+\uvec{z}\nabla^2A_z}

\begin{itemize}
\tightlist
\item
  بنابراین
\end{itemize}

\ELdisp{\nabla(\nabla\cdot\vect{A})-\nabla^2\vect{A}\ELbr =\mu_0\vect{J}}

\begin{itemize}
\tightlist
\item
  برای ساده شدن رابطه فوق انتخاب می‌کنیم
\end{itemize}

\ELdisp{\nabla\cdot\vect{A}\ELbr =0\quad\ELbr \ELbr \Longrightarrow\quad\ELbr \boxed{\nabla^2\vect{A}=-\mu_0\vect{J}}}

\begin{itemize}
\tightlist
\item
  رابطه اخیر \textbf{معادله برداری پواسون} است.
\item
  در مختصات کارتزین داریم
\end{itemize}

\ELdisp{\nabla^2A_x\ELbr =-\mu_0J_x,\qquad\ELbr \nabla^2A_y\ELbr =-\mu_0J_y,\qquad\ELbr \nabla^2A_z\ELbr =-\mu_0J_z}

\begin{itemize}
\tightlist
\item
  هر کدام از این سه معادله از نظر ریاضی شبیه معادله پواسون در الکتریسیته ساکن هستند.
\item
  پیش از این داشتیم
\end{itemize}

\ELdisp{\nabla^2V\ELbr =-\frac{\rho}{\epsilon_0}\quad\ELbr \ELbr \Longrightarrow\quad\ELbr  V\ELbr =\frac{1}{4\pi\epsilon_0}\int_{V'}\frac{\rho}{\lvert\vect{R}-\vect{R}'\rvert}\,dv'}

\end{ELunit}

\begin{ELjoin}

\begin{itemize}
\tightlist
\item
  به طور مشابه خواهیم داشت
\end{itemize}

\begin{ELimportant}{}

\ELdisp{\vect{A}\ELbr =\frac{\mu_0}{4\pi}\int_{V'}\frac{\vect{J}}{\lvert\vect{R}-\vect{R}'\rvert}\,dv'}

\end{ELimportant}

\end{ELjoin}

\begin{ELunit}

\begin{itemize}
\tightlist
\item
  به این ترتیب می‌توان با استفاده از رابطه فوق، \ELim{\vect{A}} را از روی چگالی جریان بدست آورد و سپس با گرفتن کرل از آن، \ELim{\vect{B}} را محاسبه نمود.
\item
  مفهوم فیزیکی پتانسیل مغناطیسی برداری
\end{itemize}

\ELdisp{\Phi\ELbr =\int_S\vect{B}\cdot d\vect{s}}

\end{ELunit}

\begin{ELjoin}

\begin{itemize}
\tightlist
\item
  \ELim{\Phi}: شار مغناطیسی با واحد وبر
\end{itemize}

\begin{ELimportant}{}

\ELdisp{\Phi\ELbr =\int_S(\nabla\times\vect{A})\cdot d\vect{s}\ELbr =\oint_C\vect{A}\cdot\dif\vect{\ell}\qquad\ELbr (\mathrm{Wb})}

\end{ELimportant}

\end{ELjoin}

\begin{ELunit}

\begin{itemize}
\tightlist
\item
  انتگرال خطی \ELim{\vect{A}} به دور هر مسیر بسته \ELim{C}، برابر با کل شار مغناطیسی گذرنده از سطح محصور شده توسط این مسیر است.
\end{itemize}

\end{ELunit}

\ELhold

\ELsection{قانون بیوساوار}{قانون بیوساوار}

\begin{ELunit}

\ELfigure{m05-map-biot}{}

\begin{itemize}
\tightlist
\item
  در یک سیم نازک حامل جریان \ELim{I} و با سطح مقطع \ELim{S} داریم
\end{itemize}

\ELdisp{\vect{J}\,dv'\ELbr =JS\,d\ell'\ELbr =I\,\dif\vect{\ell}'\quad\ELbr \ELbr \Longrightarrow\quad\ELbr \vect{A}\ELbr =\frac{\mu_0I}{4\pi}\oint_{C'}\frac{\dif\vect{\ell}'}{\lvert\vect{R}-\vect{R}'\rvert}}

\begin{itemize}
\tightlist
\item
  می‌توان نشان داد
\end{itemize}

\ELdisp{\vect{B}\ELbr =\nabla\times\vect{A}\ELbr =\nabla\times\left[\frac{\mu_0I}{4\pi}\oint_{C'}\frac{\dif\vect{\ell}'}{\lvert\vect{R}-\vect{R}'\rvert}\right]}

\end{ELunit}

\begin{ELimportant}{قانون بیوساوار}

\ELdisp{\vect{B}\ELbr =\frac{\mu_0I}{4\pi}\oint_{C'}\frac{\dif\vect{\ell}'\times(\vect{R}-\vect{R}')}{\lvert\vect{R}-\vect{R}'\rvert^3}}

\end{ELimportant}

\begin{ELunit}

\begin{itemize}
\tightlist
\item
  معادله فوق را \textbf{قانون بیوساوار} گویند.

  \begin{itemize}
  \tightlist
  \item
    به طور کلی استفاده از قانون بیوساوار مشکل‌تر از قانون مداری آمپر است. اما اگر مسیر بسته‌ای که \ELim{\vect{B}} روی آن ثابت باشد پیدا نشود، استفاده از قانون مداری آمپر برای تعیین \ELim{\vect{B}} امکان‌پذیر نیست.
  \end{itemize}
\end{itemize}

\end{ELunit}

\ELnextnum{5-4}

\begin{ELexample}{}

جریان مستقیم \ELim{I} در یک سیم مستقیم به طول \ELim{2L} جاری است. چگالی شار مغناطیسی را در نقطه‌ای به فاصله \ELim{r} از سیم در صفحه عمود منصف آن تعیین کنید.

\ELfigure{m05-ex4}{}

\begin{ELsolution}

\begin{itemize}
\tightlist
\item
  جریان فقط در مدارهای بسته وجود دارد. بنابراین سیم موجود در این مثال باید بخشی از یک حلقه حامل جریان باشد. چون از بقیه مدار اطلاعی نداریم، بنابراین نمی‌توان از قانون مداری آمپر استفاده کرد.
\item
  \textbf{روش اول حل}
\end{itemize}

\ELdisp{\left.\begin{aligned}&\vect{A}=\frac{\mu_0I}{4\pi}\oint_{C'}\frac{\dif\vect{\ell}'}{\lvert\vect{R}-\vect{R}'\rvert}\\&\dif\vect{\ell}'=dz'\,\uvec{z}\\&\vect{R}=r\uvec{r}\\&\vect{R}'=z'\uvec{z}\end{aligned}\right\}\quad\ELbr \begin{aligned}\vect{A}&=\uvec{z}\frac{\mu_0I}{4\pi}\int_{-L}^{L}\frac{dz'}{\sqrt{z'^2+r^2}}\\&=\uvec{z}\frac{\mu_0I}{4\pi}\ln\frac{\sqrt{L^2+r^2}+L}{\sqrt{L^2+r^2}-L}\end{aligned}}
\ELdisp{\vect{B}\ELbr =\nabla\times\vect{A}\ELbr =\nabla\times(\uvec{z}A_z)\ELbr =\uvec{r}\frac1r\frac{\partial A_z}{\partial\phi}-\uvec{\phi}\frac{\partial A_z}{\partial r}}
\ELdisp{\begin{aligned}\vect{B}&=-\uvec{\phi}\frac{\partial}{\partial r}\left[\frac{\mu_0I}{4\pi}\ln\frac{\sqrt{L^2+r^2}+L}{\sqrt{L^2+r^2}-L}\right]\\&=\uvec{\phi}\frac{\mu_0IL}{2\pi r\sqrt{L^2+r^2}}\end{aligned}}

\begin{itemize}
\tightlist
\item
  \textbf{روش دوم حل}
\end{itemize}

\ELdisp{\left.\begin{aligned}&\vect{B}=\frac{\mu_0I}{4\pi}\oint_{C'}\frac{\dif\vect{\ell}'\times(\vect{R}-\vect{R}')}{\lvert\vect{R}-\vect{R}'\rvert^3}\\&\vect{R}-\vect{R}'=r\uvec{r}-z'\uvec{z}\\&\dif\vect{\ell}'\times(\vect{R}-\vect{R}')=dz'\uvec{z}\times(r\uvec{r}-z'\uvec{z})=r\,dz'\,\uvec{\phi}\end{aligned}\right\}\quad\ELbr \begin{aligned}\vect{B}=\int d\vect{B}&=\uvec{\phi}\frac{\mu_0I}{4\pi}\int_{-L}^{L}\frac{r\,dz'}{(z'^2+r^2)^{3/2}}\\&=\uvec{\phi}\frac{\mu_0IL}{2\pi r\sqrt{L^2+r^2}}\end{aligned}}

\end{ELsolution}

\end{ELexample}

\ELnextnum{5-5}

\begin{ELexample}{}

چگالی شار مغناطیسی در مرکز یک حلقه مربعی به ضلع \ELim{w} و حامل جریان \ELim{I} را بدست آورید.

\ELfigure{m05-ex5}{}

\begin{ELsolution}

\begin{itemize}
\tightlist
\item
  فرض می‌کنیم که حلقه در صفحه \ELim{xy} باشد. \ELim{\vect{B}} در مرکز حلقه، چهار برابر چگالی شار مغناطیسی ناشی از یک ضلع به طول \ELim{w} است. بنابراین
\item
  اندازه چگالی شار مغناطیسی یک تکه سیم به طول \ELim{2L} در فاصله \ELim{r} از آن
\end{itemize}

\ELdisp{B\ELbr =\frac{\mu_0IL}{2\pi r\sqrt{L^2+r^2}}\qquad\ELbr  L\ELbr =\frac w2,\quad\ELbr  r\ELbr =\frac w2}
\ELdisp{\vect{B}\ELbr =\uvec{z}\frac{\mu_0I}{\sqrt2\pi w}\times4\ELbr =\uvec{z}\frac{2\sqrt2\mu_0I}{\pi w}}

\end{ELsolution}

\end{ELexample}

\ELnextnum{5-6}

\begin{ELexample}{}

چگالی شار مغناطیسی روی محور یک حلقه جریان دایروی به شعاع \ELim{b} و حامل جریان \ELim{I} را بدست آورید.

\ELfigure{m05-ex6}{}

\begin{ELsolution}

\ELdisp{\left.\begin{aligned}&\vect{B}=\frac{\mu_0I}{4\pi}\oint_{C'}\frac{\dif\vect{\ell}'\times(\vect{R}-\vect{R}')}{\lvert\vect{R}-\vect{R}'\rvert^3}\\&\dif\vect{\ell}'=b\,d\phi'\,\uvec{\phi'},\quad\vect{R}=z\uvec{z},\quad\vect{R}'=b\uvec{r'}\\&\vect{R}-\vect{R}'=z\uvec{z}-b\uvec{r'}\\&\begin{aligned}\dif\vect{\ell}'\times(\vect{R}-\vect{R}')&=b\,d\phi'\,\uvec{\phi'}\times(z\uvec{z}-b\uvec{r'})\\&=bz\,d\phi'\,\uvec{r'}+b^2d\phi'\,\uvec{z}\end{aligned}\end{aligned}\right\}\quad\ELbr \vect{B}\ELbr =\frac{\mu_0I}{4\pi}\int_0^{2\pi}\uvec{z}\frac{b^2\,d\phi'}{(z^2+b^2)^{3/2}}}

\begin{itemize}
\tightlist
\item
  به صورت شهودی واضح است که مؤلفه \ELim{r} چگالی شار مغناطیسی صفر است.
\end{itemize}

\begin{ELimportant}{}

\ELdisp{\vect{B}\ELbr =\uvec{z}\frac{\mu_0Ib^2}{2(z^2+b^2)^{3/2}}\qquad\ELbr (\mathrm{T})}

\end{ELimportant}

\end{ELsolution}

\end{ELexample}

\ELnextnum{5-7}

\begin{ELexample}{}

چگالی شار مغناطیسی ناشی از یک دوقطبی مغناطیسی به شعاع \ELim{b} و حامل جریان \ELim{I} را به فاصله بسیار دور از آن بدست آورید.

\ELfigure{m05-ex7}{}

\begin{ELsolution}

\ELdisp{\vect{A}\ELbr =\frac{\mu_0I}{4\pi}\oint_{C'}\frac{\dif\vect{\ell}'}{\lvert\vect{R}-\vect{R}'\rvert}}

\begin{itemize}
\tightlist
\item
  با توجه به تقارن هندسه، میدان مغناطیسی مستقل از زاویه \ELim{\phi} است. بنابراین جهت سادگی نقطه مشاهده را در صفحه \ELim{yz} قرار می‌دهیم.
\end{itemize}

\ELdisp{\dif\vect{\ell}'\ELbr =b\,d\phi'\,\uvec{\phi'}\ELbr =b\,d\phi'\left(-\uvec{x}\sin\phi'+\uvec{y}\cos\phi'\right)}
\ELdisp{\vect{A}\ELbr =-\uvec{x}\frac{\mu_0I}{4\pi}\int_0^{2\pi}\frac{b\sin\phi'}{\lvert\vect{R}-\vect{R}'\rvert}\,d\phi'+\uvec{y}\frac{\mu_0I}{4\pi}\int_0^{2\pi}\frac{b\cos\phi'}{\lvert\vect{R}-\vect{R}'\rvert}\,d\phi'}

\begin{itemize}
\tightlist
\item
  به صورت شهودی روشن است که مؤلفه \ELim{\vect{A}} در راستای \ELim{y} صفر است.
\end{itemize}

\ELfigure{m05-ex7-top}{}

\ELdisp{\begin{aligned}\vect{R}=R\uvec{R}&=R\left(\sin\theta\cos(\pi/2)\uvec{x}+\sin\theta\sin(\pi/2)\uvec{y}+\cos\theta\,\uvec{z}\right)\\&=R\left(\sin\theta\,\uvec{y}+\cos\theta\,\uvec{z}\right)\end{aligned}}
\ELdisp{\vect{R}'\ELbr =b\uvec{r'}\ELbr =b\left(\cos\phi'\,\uvec{x}+\sin\phi'\,\uvec{y}\right)}
\ELdisp{\vect{R}-\vect{R}'\ELbr =-b\cos\phi'\,\uvec{x}+(R\sin\theta-b\sin\phi')\uvec{y}+R\cos\theta\,\uvec{z}}
\ELdisp{\begin{aligned}\lvert\vect{R}-\vect{R}'\rvert^2&=R^2+b^2-2bR\sin\theta\sin\phi'\\&=R^2\left(1+\frac{b^2}{R^2}-\frac{2b}{R}\sin\theta\sin\phi'\right)\end{aligned}}
\ELdisp{\begin{aligned}\frac{1}{\lvert\vect{R}-\vect{R}'\rvert}&=\frac1R\left(1+\frac{b^2}{R^2}-\frac{2b}{R}\sin\theta\sin\phi'\right)^{-1/2}\\&\cong\frac1R\left(1-\frac{2b}{R}\sin\theta\sin\phi'\right)^{-1/2}\cong\frac1R\left(1+\frac bR\sin\theta\sin\phi'\right)\end{aligned}}
\ELdisp{\begin{aligned}\vect{A}&=-\uvec{x}\frac{\mu_0Ib}{4\pi R}\int_0^{2\pi}\left(1+\frac bR\sin\theta\sin\phi'\right)\sin\phi'\,d\phi'\\&=-\uvec{x}\frac{\mu_0Ib^2}{4R^2}\sin\theta\end{aligned}}

\begin{itemize}
\tightlist
\item
  حال در هر نقطه مشاهده دلخواهی بردار پتانسل مغناطیسی به صورت زیر خواهد شد
\end{itemize}

\ELdisp{\vect{A}\ELbr =\uvec{\phi}\frac{\mu_0Ib^2}{4R^2}\sin\theta\quad\ELbr \ELbr \Longrightarrow\quad\ELbr \vect{A}\ELbr =\uvec{\phi}\frac{\mu_0(I\pi b^2)}{4\pi R^2}\sin\theta}

\begin{itemize}
\tightlist
\item
  بردار گشتاور دوقطبی مغناطیسی
\end{itemize}

\ELdisp{\vect{m}\ELbr =\uvec{z}I\pi b^2\ELbr =\uvec{z}IS\ELbr =\uvec{z}m}

\begin{ELimportant}{}

\ELdisp{\vect{A}\ELbr =\frac{\mu_0\vect{m}\times\uvec{R}}{4\pi R^2}}

\end{ELimportant}

\ELdisp{\vect{B}\ELbr =\nabla\times\vect{A}\quad\ELbr \ELbr \Longrightarrow\quad\ELbr \vect{B}\ELbr =\frac{\mu_0Ib^2}{4R^3}\left(\uvec{R}2\cos\theta+\uvec{\theta}\sin\theta\right)}

\begin{ELimportant}{}

\ELdisp{\vect{B}\ELbr =\frac{\mu_0m}{4\pi R^3}\left(\uvec{R}2\cos\theta+\uvec{\theta}\sin\theta\right)}

\end{ELimportant}

\ELfigure{m05-dipoles}{}

\end{ELsolution}

\end{ELexample}

\ELhold

\ELsection{رفتار میدان مغناطیسی ساکن در محیط‌های مادی}{رفتار میدان مغناطیسی ساکن در محیط‌های مادی}

\begin{ELunit}

\ELfigure{m05-map-media}{}

\begin{itemize}
\tightlist
\item
  تمام مواد از اتم‌هایی با هسته با بار مثبت و تعدادی الکترون با بار منفی در حال گردش به دور آن تشکیل می‌شوند.
\item
  حرکت الکترون‌ها باعث تولید جریان گردان و در نتیجه دوقطبی‌های مغناطیسی میکروسکوپی می‌شود.
\item
  در غیاب میدان مغناطیسی خارجی، دو قطبی‌های مغناطیسی اتم‌های اکثر مواد (به جز آهنربای دائمی) دارای جهت‌های تصادفی هستند و هیچ گشتاور مغناطیسی خالصی ایجاد نمی‌شود. با اعمال میدان مغناطیسی خارجی، گشتاور مغناطیسی الکترون‌های چرخان هم‌امتداد می‌شوند.
\item
  مواد مختلف در معرض میدان مغناطیسی چه رفتاری از خود نشان می‌دهند و چگونه می‌توان این رفتار را تحلیل کرد؟
\end{itemize}

\end{ELunit}

\ELhold

\ELsubsection{مغناطیس‌شدگی و چگالی جریان‌های معادل}{مغناطیس‌شدگی و چگالی جریان‌های معادل}

\begin{ELunit}

\ELfigure{m05-magnetization}{}

\begin{itemize}
\tightlist
\item
  بردار مغناطیس شدگی (و بردار قطبی شدگی در الکتریسیته ساکن)
\end{itemize}

\ELdisp{\vect{M}\ELbr =\lim_{\Delta v\to0}\frac{\sum_{k=1}^{n\Delta v}\vect{m}_k}{\Delta v}\quad\ELbr (\mathrm{A/m})\qquad\ELbr \qquad\ELbr \vect{P}\ELbr =\lim_{\Delta v\to0}\frac{\sum_{k=1}^{n\Delta v}\vect{p}_k}{\Delta v}\quad\ELbr (\mathrm{C/m^2})}

\begin{itemize}
\tightlist
\item
  پتانسیل مغناطیسی برداری ناشی از ماده مغناطیس شده (و پتانسیل الکتریکی ناشی از ماده قطبی شده)
\end{itemize}

\ELdisp{\vect{A}\ELbr =\frac{\mu_0}{4\pi}\int_{V'}\frac{\nabla'\times\vect{M}}{R}\,dv'+\frac{\mu_0}{4\pi}\oint_{S'}\frac{\vect{M}\times\uvec{n}'}{R}\,ds'}
\ELdisp{V\ELbr =\frac{1}{4\pi\epsilon_0}\oint_{S'}\frac{\vect{P}\cdot\uvec{n}'}{R}\,ds'+\frac{1}{4\pi\epsilon_0}\int_{V'}\frac{(-\nabla'\cdot\vect{P})}{R}\,dv'}

\end{ELunit}

\begin{ELjoin}

\begin{itemize}
\tightlist
\item
  چگالی جریان‌های سطحی و حجمی مغناطیس‌شدگی معادل
\end{itemize}

\begin{ELimportant}{}

\ELdisp{\vect{J}_{ms}\ELbr =\vect{M}\times\uvec{n},\qquad\ELbr \vect{J}_m\ELbr =\nabla\times\vect{M}}

\end{ELimportant}

\end{ELjoin}

\begin{ELunit}

\begin{itemize}
\tightlist
\item
  چگالی بارهای سطحی و حجمی قطبی‌شدگی معادل
\end{itemize}

\ELdisp{\rho_{ps}\ELbr =\vect{P}\cdot\uvec{n},\qquad\ELbr \rho_p\ELbr =-\nabla\cdot\vect{P}}

\begin{itemize}
\tightlist
\item
  مسئله یافتن چگالی شار مغناطیسی \ELim{\vect{B}} ناشی از یک چگالی حجمی مشخص از گشتاور دو قطبی مغناطیسی \ELim{\vect{M}}، به پیدا کردن چگالی‌های جریان مغناطیس‌شدگی معادل و سپس تعیین \ELim{\vect{A}} و به دنبال آن محاسبه \ELim{\vect{B}} تبدیل می‌شود.
\end{itemize}

\end{ELunit}

\ELnextnum{5-8}

\begin{ELexample}{}

چگالی شار مغناطیسی را روی محور یک استوانه مغناطیسی بدست آورید. این استوانه دارای شعاع \ELim{b}، طول \ELim{L} و بردار مغناطیس شدگی \ELim{\vect{M}=\uvec{z}M_0} است.

\ELfigure{m05-ex8}{}

\begin{ELsolution}

\ELdisp{\vect{J}_m\ELbr =\nabla\times\vect{M}\ELbr =0}

\begin{itemize}
\tightlist
\item
  چگالی جریان سطحی معادل روی دیواره جانبی
\end{itemize}

\ELdisp{\vect{J}_{ms}\ELbr =\vect{M}\times\uvec{n}\ELbr =\uvec{z}M_0\times\uvec{r}\ELbr =\uvec{\phi}M_0}

\begin{itemize}
\tightlist
\item
  چگالی جریان سطحی معادل روی سطوح بالایی و پایینی صفر است.
\end{itemize}

\ELdisp{\left.\begin{aligned}&\vect{B}=\frac{\mu_0}{4\pi}\int_{S'}\frac{\vect{J}_{ms}\times(\vect{R}-\vect{R}')\,ds'}{\lvert\vect{R}-\vect{R}'\rvert^3}\\&ds'=b\,d\phi'\,dz'\\&\vect{R}=z\uvec{z}\\&\vect{R}'=b\uvec{r'}+z'\uvec{z}\\&\vect{R}-\vect{R}'=(z-z')\uvec{z}-b\uvec{r'}\\&\begin{aligned}\vect{J}_{ms}\times(\vect{R}-\vect{R}')&=M_0\uvec{\phi'}\times\left((z-z')\uvec{z}-b\uvec{r'}\right)\\&=M_0(z-z')\uvec{r'}+bM_0\uvec{z}\end{aligned}\end{aligned}\right\}}

\begin{itemize}
\tightlist
\item
  به صورت شهودی واضح است که مؤلفه \ELim{r} چگالی شار مغناطیسی صفر است.
\end{itemize}

\ELdisp{\begin{aligned}\vect{B}=\int d\vect{B}&=\uvec{z}\int_0^L\frac{\mu_0M_0b^2\,dz'}{2\left[(z-z')^2+b^2\right]^{3/2}}\\&=\uvec{z}\frac{\mu_0M_0}{2}\left[\frac{z}{\sqrt{z^2+b^2}}-\frac{z-L}{\sqrt{(z-L)^2+b^2}}\right]\end{aligned}}

\end{ELsolution}

\end{ELexample}

\ELhold

\ELsection{شدت میدان مغناطیسی و نفوذپذیری نسبی}{شدت میدان مغناطیسی و نفوذپذیری نسبی}

\begin{ELunit}

\ELfigure{m05-map-hmu}{}

\end{ELunit}

\ELhold

\ELsubsection{شدت میدان مغناطیسی}{شدت میدان مغناطیسی}

\begin{ELunit}

\begin{itemize}
\tightlist
\item
  چون اعمال یک میدان مغناطیسی خارجی باعث هم‌امتداد شدن گشتاورهای دوقطبی داخلی و القای گشتاور مغناطیسی در ماده مغناطیسی می‌شود، انتظار داریم که چگالی شار مغناطیسی حاصل شده در حضور ماده مغناطیسی با مقدار آن در فضای آزاد متفاوت باشد.
\item
  با درنظر گرفتن چگالی جریان حجمی معادل در ماده داریم
\end{itemize}

\ELdisp{\frac{1}{\mu_0}\nabla\times\vect{B}\ELbr =\vect{J}+\vect{J}_m\ELbr =\vect{J}+\nabla\times\vect{M}\quad\ELbr \ELbr \Longrightarrow\quad\ELbr \nabla\times\left(\frac{\vect{B}}{\mu_0}-\vect{M}\right)\ELbr =\vect{J}}

\end{ELunit}

\begin{ELjoin}

\begin{itemize}
\tightlist
\item
  شدت میدان مغناطیسی را به صورت زیر تعریف می‌کنیم
\end{itemize}

\begin{ELimportant}{}

\ELdisp{\vect{H}\ELbr =\frac{\vect{B}}{\mu_0}-\vect{M}\qquad\ELbr (\mathrm{A/m})}

\end{ELimportant}

\end{ELjoin}

\begin{ELunit}

\begin{itemize}
\tightlist
\item
  استفاده از بردار \ELim{\vect{H}} ما را قادر می‌سازد که معادله کرل ارتباط دهنده میدان مغناطیسی و توزیع جریان‌های آزاد را در هر محیطی بدون نیاز به بکار بردن \ELim{\vect{M}} یا \ELim{\vect{J}_m} بنویسیم
\end{itemize}

\end{ELunit}

\begin{ELimportant}{}

\ELdisp{\nabla\times\vect{H}\ELbr =\vect{J}\qquad\ELbr (\mathrm{A/m^2})}

\end{ELimportant}

\begin{ELjoin}

\begin{itemize}
\tightlist
\item
  \ELim{\vect{J}}: چگالی جریان آزاد
\item
  فرم انتگرالی
\end{itemize}

\ELdisp{\int_S(\nabla\times\vect{H})\cdot d\vect{s}\ELbr =\int_S\vect{J}\cdot d\vect{s}\quad\ELbr \ELbr \Longrightarrow\quad\ELbr \boxed{\oint_C\vect{H}\cdot\dif\vect{\ell}=I\qquad(\mathrm{A})}}

\begin{itemize}
\tightlist
\item
  رابطه اخیر شکل دیگر قانون مداری آمپر است.
\item
  معادلات اصلی حاکم بر مغناطیس ساکن در هر محیطی
\end{itemize}

\begin{ELimportant}{}

\ELdisp{\left\{\begin{aligned}&\nabla\cdot\vect{B}=0\\&\nabla\times\vect{H}=\vect{J}\end{aligned}\right.}

\end{ELimportant}

\end{ELjoin}

\ELhold

\ELsubsection{ضریب نفوذپذیری نسبی}{ضریب نفوذپذیری نسبی}

\begin{ELjoin}

\begin{itemize}
\tightlist
\item
  رابطه بین بردار مغناطیس شدگی و شدت میدان مغناطیسی
\end{itemize}

\ELdisp{\vect{M}\ELbr =\chi_m\vect{H}}

\begin{itemize}
\tightlist
\item
  \ELim{\chi_m}: ضریب حساسیت مغناطیسی
\item
  رابطه بین شدت میدان مغناطیسی و چگالی شار مغناطیسی
\end{itemize}

\begin{ELimportant}{}

\ELdisp{\begin{aligned}\vect{B}&=\mu_0(1+\chi_m)\vect{H}\\&=\mu_0\mu_r\vect{H}=\mu\vect{H}\qquad(\mathrm{Wb/m^2})\end{aligned}}

\end{ELimportant}

\end{ELjoin}

\begin{ELunit}

\ELdisp{\mu_r\ELbr =1+\chi_m\ELbr =\frac{\mu}{\mu_0}}

\begin{itemize}
\tightlist
\item
  \ELim{\mu_r}: کمیت بدون بعد به نام نفوذپذیری نسبی محیط
\item
  روابط میدان الکتریکی و مغناطیسی ساکن دوگان یکدیگرند
\end{itemize}

\end{ELunit}

\Needspace{19\baselineskip}
\begin{ELltr}
\begin{longtable}[c]{cc}
\toprule\noalign{}
\ELtablehead \lr{Electrostatics} & \lr{Magnetostatics} \\\noalign{\vskip 4pt}
\midrule\noalign{}
\endhead
\bottomrule\noalign{}
\endlastfoot
\ELim{\vect{E}} & \ELim{\vect{B}} \\\noalign{\vskip 4pt}
\ELim{\vect{D}} & \ELim{\vect{H}} \\\noalign{\vskip 4pt}
\ELim{\epsilon} & \ELim{\dfrac1\mu} \\\noalign{\vskip 4pt}
\ELim{\vect{P}} & \ELim{-\vect{M}} \\\noalign{\vskip 4pt}
\ELim{\rho} & \ELim{\vect{J}} \\\noalign{\vskip 4pt}
\ELim{V} & \ELim{\vect{A}} \\\noalign{\vskip 4pt}
\ELim{\cdot} & \ELim{\times} \\\noalign{\vskip 4pt}
\ELim{\times} & \ELim{\cdot} \\\noalign{\vskip 4pt}
\end{longtable}\end{ELltr}


\ELhold

\ELsubsection{رفتار مواد مغناطیسی}{رفتار مواد مغناطیسی}

\begin{ELunit}

\begin{itemize}
\tightlist
\item
  مواد مغناطیسی را به طور تقریبی می‌توان بر اساس مقدار ضریب نفوذپذیری نسبی آن‌ها به سه گروه اصلی طبقه بندی کرد

  \begin{itemize}
  \tightlist
  \item
    \textbf{مواد دیامغناطیس}: \ELim{\mu_r<1} (\ELim{\chi_m} عدد منفی بسیار کوچکی است)

    \begin{itemize}
    \tightlist
    \item
      در این مواد گشتاور مغناطیسی خالص در غیاب میدان مغناطیسی اعمال شده خارجی صفر است. \ELim{\ELto} مس، سرب، جیوه، نقره و طلا \ELim{\ELto} \ELim{\chi_m\approx-10^{-5}}
    \end{itemize}
  \item
    \textbf{مواد پارامغناطیس}: \ELim{\mu_r>1} (\ELim{\chi_m} عدد مثبت بسیار کوچکی است)

    \begin{itemize}
    \tightlist
    \item
      در این مواد گشتاور مغناطیسی متوسط خالصی در غیاب میدان مغناطیسی اعمالی خارجی وجود دارد. \ELim{\ELto} آلومینیوم، منیزیوم و تنگستن \ELim{\ELto} \ELim{\chi_m\approx10^{-5}}
    \end{itemize}
  \item
    \textbf{مواد فرومغناطیس}: \ELim{\mu_r\gg1} (\ELim{\chi_m} عدد مثبت بزرگی است)

    \begin{itemize}
    \tightlist
    \item
      در این مواد گشتاور مغناطیسی از نظر اندازه چندین برابر مقدار نظیر در مواد پارامغناطیس است. \ELim{\ELto} کوبالت، نیکل و آهن
    \item
      رابطه بین \ELim{\vect{B}} و \ELim{\vect{H}} در یک ماده فرومغناطیس غیرخطی است.
    \end{itemize}
  \end{itemize}
\end{itemize}

\ELfigure{m05-hysteresis}{پدیده هیسترزیس یا پس‌ماند}

\end{ELunit}

\ELhold

\ELsection{مدارهای مغناطیسی}{مدارهای مغناطیسی}

\begin{ELunit}

\ELfigure{m05-map-circ}{}

\begin{itemize}
\tightlist
\item
  مشابه بحث مدارهای الکتریکی، گاهی اوقات لازم است که در یک مدار مغناطیسی، شار مغناطیسی و شدت میدان مغناطیسی ناشی از سیم‌پیچ‌های حامل جریان محاسبه شود.
\item
  برای تحلیل مدارهای مغناطیسی داریم
\end{itemize}

\ELdisp{\nabla\cdot\vect{B}\ELbr =0,\qquad\ELbr \nabla\times\vect{H}\ELbr =\vect{J}}
\ELdisp{\oint_C\vect{H}\cdot\dif\vect{\ell}\ELbr =NI\ELbr =\mathcal{V}_m}

\end{ELunit}

\begin{ELjoin}

\begin{itemize}
\tightlist
\item
  \ELim{\mathcal{V}_m}: نیروی محرکه مغناطیسی با واحد \ELim{\mathrm{A}}. این کمیت مشابه نیروی محرکه الکتریکی در مدارات الکتریکی است.
\end{itemize}

\ELnextnum{5-9}

\begin{ELexample}{}

فرض کنید که \ELim{N} دور سیم برروی یک هسته چنبره‌ای از ماده مغناطیسی با نفوذپذیری \ELim{\mu} پیچیده شده باشد. هسته دارای شعاع متوسط \ELim{r_0}، سطح مقطع دایروی با شعاع \ELim{a} و یک فاصله هوایی کوچک به طول \ELim{\ell_g} است. جریان دائمی \ELim{I_0} در سیم جاری است. مطلوبست

\begin{itemize}
\tightlist
\item
  الف. چگالی شار مغناطیسی در هسته مغناطیسی
\item
  ب. شدت میدان مغناطیسی در هسته مغناطیسی
\item
  ج. شدت میدان مغناطیسی در فاصله هوایی
\end{itemize}

\ELfigure{m05-ex9}{}

\begin{ELsolution}

\begin{itemize}
\tightlist
\item
  \textbf{فرض اول}: از شار نشتی صرف نظر می‌کنیم
\end{itemize}

\ELfigure{m05-ex9-gauss}{}

\ELdisp{\nabla\cdot\vect{B}\ELbr =0\quad\ELbr \ELbr \Longrightarrow\quad\ELbr \oint_S\vect{B}\cdot d\vect{s}\ELbr =0}

\begin{itemize}
\tightlist
\item
  شار در هسته و فاصله هوایی یکسان است.
\item
  \textbf{فرض دوم}: از اثرات حاشیه‌ای شار در فاصله هوایی صرف نظر می‌کنیم.

  \begin{itemize}
  \tightlist
  \item
    چگالی شار در هسته و فاصله هوایی یکسان است.
  \end{itemize}
\end{itemize}

\ELdisp{\vect{B}_f\ELbr =\vect{B}_g\ELbr =\uvec{\phi}B_f}
\ELdisp{\vect{H}_f\ELbr =\uvec{\phi}\frac{B_f}{\mu},\qquad\ELbr \vect{H}_g\ELbr =\uvec{\phi}\frac{B_f}{\mu_0}}
\ELdisp{\oint_C\vect{H}\cdot\dif\vect{\ell}\ELbr =NI_0\quad\ELbr \ELbr \Longrightarrow\quad\ELbr \frac{B_f}{\mu}(2\pi r_0-\ell_g)+\frac{B_f}{\mu_0}\ell_g\ELbr =NI_0}

\begin{itemize}
\tightlist
\item
  الف.
\end{itemize}

\ELdisp{\vect{B}_f\ELbr =\uvec{\phi}\frac{\mu_0\mu NI_0}{\mu_0(2\pi r_0-\ell_g)+\mu\ell_g}}

\begin{itemize}
\tightlist
\item
  ب.
\end{itemize}

\ELdisp{\vect{H}_f\ELbr =\uvec{\phi}\frac{\mu_0NI_0}{\mu_0(2\pi r_0-\ell_g)+\mu\ell_g}}

\begin{itemize}
\tightlist
\item
  ج.
\end{itemize}

\ELdisp{\vect{H}_g\ELbr =\uvec{\phi}\frac{\mu NI_0}{\mu_0(2\pi r_0-\ell_g)+\mu\ell_g}}

\end{ELsolution}

\end{ELexample}

\end{ELjoin}

\begin{ELjoin}

\begin{itemize}
\tightlist
\item
  اگر شعاع سطح مقطع هسته بسیار کوچکتر از شعاع متوسط چنبره باشد، چگالی شار مغناطیسی در هسته تقریباً ثابت بوده و داریم
\end{itemize}

\ELdisp{\Phi\ELbr \cong BS\quad\ELbr \ELbr \Longrightarrow\quad\ELbr \Phi\ELbr =\frac{NI_0}{(2\pi r_0-\ell_g)/\mu S+\ell_g/\mu_0S}\quad\ELbr \ELbr \Longrightarrow\quad\ELbr \Phi\ELbr =\frac{\mathcal{V}_m}{\mathcal{R}_f+\mathcal{R}_g}}

\begin{ELimportant}{}

\ELdisp{\mathcal{R}_f\ELbr =\frac{2\pi r_0-\ell_g}{\mu S}\ELbr =\frac{\ell_f}{\mu S},\qquad\ELbr \mathcal{R}_g\ELbr =\frac{\ell_g}{\mu_0S}}

\end{ELimportant}

\end{ELjoin}

\begin{ELunit}

\begin{itemize}
\item
  \ELim{\mathcal{R}_f} و \ELim{\mathcal{R}_g}: رلوکتانس با واحد \ELim{1/\mathrm{H}}
\item
  تشابه این مدار مغناطیسی با یک مدار الکتریکی به صورت زیر است
\end{itemize}

\ELfigure{m05-circuits}{}

\ELdisp{\Phi\ELbr =\frac{\mathcal{V}_m}{\mathcal{R}_f+\mathcal{R}_g},\qquad\ELbr  I\ELbr =\frac{\mathcal{V}}{R_f+R_g}}

\end{ELunit}

\Needspace{12\baselineskip}
\begin{ELltr}
\begin{longtable}[c]{ll}
\toprule\noalign{}
\ELtablehead \lr{Magnetic Circuits} & \lr{Electric Circuits} \\\noalign{\vskip 4pt}
\midrule\noalign{}
\endhead
\bottomrule\noalign{}
\endlastfoot
\lr{mmf}, \ELim{\mathcal{V}_m\,(=NI)} & \lr{emf}, \ELim{\mathcal{V}} \\\noalign{\vskip 4pt}
\lr{magnetic flux}, \ELim{\Phi} & \lr{electric current}, \ELim{I} \\\noalign{\vskip 4pt}
\lr{reluctance}, \ELim{\mathcal{R}} & \lr{resistance}, \ELim{R} \\\noalign{\vskip 4pt}
\lr{permeability}, \ELim{\mu} & \lr{conductivity}, \ELim{\sigma} \\\noalign{\vskip 4pt}
\end{longtable}\end{ELltr}


\begin{ELjoin}

\begin{itemize}
\tightlist
\item
  مشابه قانون ولتاژ کرشهف در مدارات الکتریکی، در مدارات مغناطیسی داریم
\end{itemize}

\begin{ELimportant}{}

\ELdisp{\sum_jN_jI_j\ELbr =\sum_k\mathcal{R}_k\Phi_k}

\end{ELimportant}

\end{ELjoin}

\begin{ELunit}

\begin{itemize}
\tightlist
\item
  در نتیجه جمع جبری آمپردورها در پیرامون هر مسیر بسته یک مدار مغناطیسی، برابر با جمع جبری حاصلضرب رلوکتانس‌ها و شارهاست.
\item
  مشابه قانون جریان کرشهف در مدارات الکتریکی، در مدارات مغناطیسی داریم
\end{itemize}

\end{ELunit}

\begin{ELimportant}{}

\ELdisp{\sum_j\Phi_j\ELbr =0}

\end{ELimportant}

\begin{ELunit}

\begin{itemize}
\tightlist
\item
  که بیان می‌کند جمع جبری شارهای مغناطیسی خارج‌شونده از یک گره در هر مدار مغناطیسی صفر است.
\item
  با وجود این تشابه ساده، تحلیل دقیق مدارهای مغناطیسی به دلایل زیر بسیار دشوار است

  \begin{itemize}
  \tightlist
  \item
    وجود شارهای نشتی
  \item
    وجود اثرات لبه‌ای که باعث می‌شود خطوط شار مغناطیسی در شکاف هوایی پخش و برآمده شود.
  \item
    وابسته بودن نفوذپذیری مواد فرومغناطیسی به شدت میدان مغناطیسی و به عبارتی غیرخطی بودن رابطه \ELim{\vect{B}} و \ELim{\vect{H}}
  \end{itemize}
\end{itemize}

\end{ELunit}

\ELnextnum{5-10}

\begin{ELexample}{}

مدار مغناطیسی شکل زیر را در نظر بگیرید. جریان‌های دائمی \ELim{I_1} و \ELim{I_2} به ترتیب در سیم پیچ‌هایی با \ELim{N_1} و \ELim{N_2} دور جاری هستند. هسته دارای سطح مقطعی به مساحت \ELim{S_c} و نفوذپذیری \ELim{\mu} است. شار مغناطیسی را در بازوی وسط تعیین کنید.

\ELfigure{m05-ex10}{}

\begin{ELsolution}

\begin{itemize}
\tightlist
\item
  مدار معادل
\end{itemize}

\ELfigure{m05-ex10-circuit}{}

\ELdisp{\mathcal{R}_1\ELbr =\frac{\ell_1}{\mu S_c},\qquad\ELbr \mathcal{R}_2\ELbr =\frac{\ell_2}{\mu S_c},\qquad\ELbr \mathcal{R}_3\ELbr =\frac{\ell_3}{\mu S_c}}
\ELdisp{\begin{aligned}&\mathit{Loop}~1:&N_1I_1&=(\mathcal{R}_1+\mathcal{R}_3)\Phi_1+\mathcal{R}_1\Phi_2\\&\mathit{Loop}~2:&N_1I_1-N_2I_2&=\mathcal{R}_1\Phi_1+(\mathcal{R}_1+\mathcal{R}_2)\Phi_2\end{aligned}}
\ELdisp{\Phi_1\ELbr =\frac{\mathcal{R}_2N_1I_1+\mathcal{R}_1N_2I_2}{\mathcal{R}_1\mathcal{R}_2+\mathcal{R}_1\mathcal{R}_3+\mathcal{R}_2\mathcal{R}_3}}

\end{ELsolution}

\end{ELexample}

\ELhold

\ELsection{شرایط مرزی میدان‌های مغناطیسی ساکن}{شرایط مرزی میدان‌های مغناطیسی ساکن}

\begin{ELunit}

\ELfigure{m05-map-bc}{}

\begin{itemize}
\tightlist
\item
  شرط مرزی مؤلفه‌های عمودی
\end{itemize}

\ELdisp{\nabla\cdot\vect{B}\ELbr =0\quad\ELbr \ELbr \Longrightarrow\quad\ELbr \boxed{B_{1n}=B_{2n}\qquad(\mathrm{T})}}

\end{ELunit}

\begin{ELjoin}

\begin{itemize}
\tightlist
\item
  در محیط‌های خطی
\end{itemize}

\begin{ELimportant}{}

\ELdisp{\mu_1H_{1n}\ELbr =\mu_2H_{2n}}

\end{ELimportant}

\end{ELjoin}

\begin{ELunit}

\begin{itemize}
\tightlist
\item
  شرط مرزی مؤلفه‌های مماسی
\end{itemize}

\ELfigure{m05-bc}{}

\ELdisp{\oint_C\vect{H}\cdot\dif\vect{\ell}\ELbr =I}
\ELdisp{\oint_{abcda}\vect{H}\cdot\dif\vect{\ell}\ELbr =\vect{H}_1\cdot\Delta\vect{w}+\vect{H}_2\cdot(-\Delta\vect{w})\ELbr =J_{sn}\Delta w}
\ELdisp{H_{1t}-H_{2t}\ELbr =J_{sn}\qquad\ELbr (\mathrm{A/m})}

\end{ELunit}

\begin{ELimportant}{}

\ELdisp{\uvec{n2}\times(\vect{H}_1-\vect{H}_2)\ELbr =\vect{J}_s\qquad\ELbr (\mathrm{A/m})}

\end{ELimportant}
